\BourbakiExerciseGroup

\begin{enumerate}
\def\labelenumi{\arabic{enumi})}
\item
  证明：§ 9 习题 20 c) 所定义的可数豪斯多夫空间是连通的。
\item
  \begin{enumerate}
  \def\labelenumii{\alph{enumii})}
  \tightlist
  \item
    设 X 为拓扑空间，\(\mathfrak{C}\) 为 X 的拓扑。证明：若 X 在半正则拓扑 \(\mathfrak{C}^*\)（\(\mathfrak{C}\) 的相伴随拓扑； \(\S 8\) ，习题 20）下是连通的，则 X 在原来的拓扑 \(\mathfrak{C}\) 下是连通的。由此推出：存在连通的次极大空间（ \(\S 8\) ，习题 22）。
  \end{enumerate}
\end{enumerate}

*b) 特别地，在实直线 \(\mathbf{R}\) 上构造一个拓扑，使得 \(\mathbf{R}\) 在其下是连通的，但 \(\mathbf{R}\) 的任何点都没有连通邻域的基本系{[}参见第四章，§ 2，习题 14 c){]}。*

\begin{enumerate}
\def\labelenumi{\arabic{enumi})}
\setcounter{enumi}{2}
\tightlist
\item
  设 A、B 为拓扑空间 X 的两个子集。
\end{enumerate}

\begin{enumerate}
\def\labelenumi{\alph{enumi})}
\item
  证明：若 A 与 B 在 X 中都是闭的，且 A ∪ B 与 A ∩ B 都是连通的，则 A 与 B 都是连通的。举例说明：当 A、B 中有一个不是闭集时，该命题不必再成立。
\item
  设 A 与 B 都是连通的，且 \(\overline{A} \cap B \neq \varnothing\)。证明：\(A \cup B\) 是连通的。
\end{enumerate}

\begin{enumerate}
\def\labelenumi{\arabic{enumi})}
\setcounter{enumi}{3}
\tightlist
\item
  设 X 为至少有两个点的连通空间。
\end{enumerate}

\begin{enumerate}
\def\labelenumi{\alph{enumi})}
\item
  设 A 为 X 的连通子集，B 为 \(\mathbb{C}\mathbf{A}\) 的一个在 \(\mathbb{C}\mathbf{A}\) 中既开又闭的子集。证明：\(\mathbf{A} \cup \mathbf{B}\) 是连通的。（把 §3 习题 5 应用于 \(\mathbf{Y} = \mathbf{A} \cup \mathbf{B}\) 与 \(\mathbf{Z} = \mathbb{C}\mathbf{A}\)。）
\item
  设 A 为 X 的连通子集，B 为集 CA 的一个分支。证明：CB 是连通的{[}利用 a){]}。
\item
  由 b) 推出：存在 X 的两个连通子集 M、N，它们都异于 X，且满足 \(M \cup N = X\) 与 \(M \cap N = \varnothing\)。
\end{enumerate}

*5) a) 在实平面 \(\mathbf{R}^2\) 中给出连通集的一个递减序列 \((\mathbf{A}_n)\) 的例子，使其交不是连通的（参见第二章，§ 4，习题 14）。

\begin{enumerate}
\def\labelenumi{\alph{enumi})}
\setcounter{enumi}{1}
\tightlist
\item
  在 \(\mathbf{R}^2\) 中，设 \(\mathbf{X}\) 为开半平面 \(y > 0\) 与 \(y < 0\) 以及诸点 \((n, 0)\)（其中 \(n\) 是整数 \(> 0\)）的并所成之集。
\end{enumerate}

X 在拓扑 \(\mathfrak{C}_0\)（由 \(\mathbb{R}^2\) 的拓扑诱导的）下是连通的。定义 X 上的一个拓扑序列 \((\mathfrak{C}_n)\)，使得 \(\mathfrak{C}_{n+1}\) 比 \(\mathfrak{C}_n\) 细，且 X 在每个 \(\mathfrak{C}_n\) 下都连通，但在序列 \((\mathfrak{C}_n)\) 的上确界拓扑下不连通{[}从 \(\mathfrak{C}_n\) 到 \(\mathfrak{C}_{n+1}\) 的过渡，是通过引入点 \((n+1,0)]\) 的新邻域实现的。*

\begin{enumerate}
\def\labelenumi{\arabic{enumi})}
\setcounter{enumi}{5}
\item
  设 X、Y 为两个连通空间，A（相应地，B）为 X（相应地，Y）的异于 X（相应地，Y）的子集。证明：在积空间 \(X \times Y\) 中，\(A \times B\) 的余集是连通的。
\item
  设 X、Y 为两个连通空间，f 为 \(X \times Y\) 到拓扑空间 Z 中的一个映射，使得所有偏映射 \(f(0, y) : X \to Z\)、\(f(x, 0) : Y \to Z (x \in Y, y \in Y)\) 都连续。证明：\(f(X \times Y)\) 是连通的。
\item
  在 § 4 习题 9 给出的积集 \(\mathbf{X} = \prod_{i\in I}\mathbf{X}_i\) 上的拓扑的例子中，设 I 是无限的且基数 t 是不可数的。证明：若诸 \(\mathbf{X}_i\) 都是连通的、正则的，且每个都至少有两个不同的点，则 X 的点 \(a = (a_i)\) 的连通分支是由这样的 \(x = (x_i)\) 所成之集：\(x_{i} = a_{i}\) 除去对有限多个指标 \(i\in I\) 外成立。（化归到 I = N 的情形；考察 X 的两个点 \(b = (b_n)\)、\(c = (c_n)\)，使得 \(b_{n}\neq c_{n}\) 对所有 n 成立。~对每个 n，设 \((\mathrm{V}_{nm})_{m\geqslant 0}\) 是 \(b_{n}\) 在 \(\mathbf{X}_n\) 中的开邻域序列，这些邻域不含 \(c_{n}\)，且使得 \(\overline{\mathbf{V}}_{n,m + 1}\subset \mathbf{V}_{nm}\)。设 Z 为 X 的具有下述性质的点 \((x_{n})\) 所成之集：存在整数 \(k > 0\)，使得 \(x_{n}\in \mathrm{V}_{n,n - k}\) 对所有满足 \(n\geqslant k\) 的 n 成立。证明：Z 在 X 中既开又闭。）
\end{enumerate}

*9) 证明：在 § 10 习题 20 所定义的局部紧空间 \(\mathbf{X}\) 中，存在点 \(x\)，其在 \(\mathbf{X}\) 中的分支不同于 \(\mathbf{X}\) 的既开又闭且含 \(x_*\) 的子集之交。

\begin{enumerate}
\def\labelenumi{\arabic{enumi})}
\setcounter{enumi}{9}
\item
  证明：线性有序集 X 赋以拓扑 \(\mathcal{G}_{+}(X)\) 或拓扑 \(\mathcal{G}_{-}(X)\)（§ 2，习题 5）后是全不连通的。
\item
  设 X 为拓扑空间，R 为 X 上的等价关系，使得每个 mod R 等价类都包含在 X 的一个分支中。证明：X/R 的诸分支是 X 的诸分支的典范像（参见第三章，§ 2，习题 17）。
\item
  设 X 为拓扑空间。证明下列三个条件等价：\(\alpha\) ) X 的诸分支是开集；\(\beta\) ) 每个 \(x \in X\) 都有一个邻域，含于每个既开又闭且含 \(x; \gamma\) ) 对每个 \(x \in X\) ，既开又闭且含 x 的诸集之交是一个开集。
\end{enumerate}

*13) 设 \(\theta\) 为一无理数。设 \(f\) 为实直线 \(\mathbf{R}\) 到环面 \(\mathbf{T}^2\) 中的映射，像 \(f(t)\)（每个 \(t \in \mathbf{R}\) 的）是 \(\mathbf{T}^2\) 中 \((t, \theta t) \in \mathbf{R}^2\) 的典范像。\(f\) 是单射且连续的。证明：在 \(\mathbf{R}\) 上赋以 \(f\) 下 \(\mathbf{T}^2\) 的拓扑的逆像拓扑时，\(\mathbf{R}\) 是连通的，但 \(\mathbf{R}\) 的任何点都没有连通邻域的基本系。*

\begin{enumerate}
\def\labelenumi{\arabic{enumi})}
\setcounter{enumi}{13}
\tightlist
\item
  \begin{enumerate}
  \def\labelenumii{\alph{enumii})}
  \tightlist
  \item
    设 X 为局部紧且局部连通的空间，A 为 X 的闭子集。设 A* 为 A 与 CA 的相对紧的诸分支之并。若 A* = A，则称 A 为满的。证明：A* 是闭的且满的。
  \end{enumerate}
\end{enumerate}

\begin{enumerate}
\def\labelenumi{\alph{enumi})}
\setcounter{enumi}{1}
\item
  再设 X 是连通的。证明：若 A 是 X 的紧子集，则 A* 是紧的（考察 A 的一个紧邻域 V 以及与 V 的边缘相交的 CA 的分支）。
\item
  再设 X 是连通且 \(\sigma\) -紧的。证明：存在满的连通紧集的一个递增序列 \((\mathbf{K}_{n})_{n\geqslant0}\)，其并为 X，且使得 \(K_{n}\subset\mathring{K}_{n+1}\) 对所有 n 成立。
\end{enumerate}

*d) 证明：若只设 X 是连通且局部紧的，或 X 是豪斯多夫、连通且局部连通的，则 b) 的结论不再成立。{[}考察 \(\mathbf{R}^2\) 的由这样的点 \((x, y)\) 组成的子空间：或者 \(y = 0\) 且 \(0 \leqslant x \leqslant 1\)，或者 \(y > 0\) 且 \(x = 1/n\)（其中 \(n\) 是整数 \(> 0\)）。*

¶ 15) 设 X 为连通的、局部连通的空间。

\begin{enumerate}
\def\labelenumi{\alph{enumi})}
\item
  设 M、N 为 X 的两个不相交的非空闭子集。证明：\(\mathbb{C}(M \cup N)\) 有一个分支，其闭包与 M 和 N 都相交。{[}若 C 是 \(\mathbb{C}(M \cup N)\) 的一个其闭包不与 M 相交的分支，注意 C 在 CN 中既开又闭{]}。
\item
  若 A 是 X 的异于 X 的闭子集，由 a) 推出：A 的每个分支都与 \(\mathbb{C}\mathbb{A}\) 的闭包相交。
\end{enumerate}

*16) 设 X 为实平面 \(\mathbf{R}^2\) 的子空间，它是直线 \(y = 0\)、以 (0, 1) 与 \((n, 1/(1 + n))\) 为端点的闭线段、以及半直线 \(x \leqslant n\)，\(y = 1/(1 + n)\)（ \(n\) 为整数 \(>0\)）的并。证明：X 是连通空间，但 X 中 (0, 1) 的余集有一个分支 C 使得 (0, 1) \(\notin \overline{\mathbf{C}}\){[}与习题 4 a) 及 15 b) 相比较{]}。*

\(^{*}\) 17) 设 \(I_{0}\) 为 R 中的区间 \([-1,1]\)，设 \(P_{0}\) 为把 \(I_{0}\) 的两点 1/2 与 1 等同起来所得的商空间；设 \(B_{0}\) 为把 \(I_{0}\) 的点 1/2 与 1 等同且把 -1/2 与 -1 等同所得的商空间。设 I 为 R 中的开区间 {]}--- I, I{[}，并设 P（相应地，B）为 I 在 \(P_{0}\)（相应地，\(B_{0}\)）中的典范像。

\begin{enumerate}
\def\labelenumi{\alph{enumi})}
\item
  证明：三个空间 I、P、B 中任何两个都不同胚（为证 I 与 P 不同胚，注意 P 中有些点，其余集是连通的；对其他空间对亦然）。
\item
  设 X 为可数无限个同胚于 I 的空间与可数无限个同胚于 B 的空间的和。设 Y 为 X 与 P 的和。证明：X 与 Y 不同胚，但 (i) 存在 X 到 Y 上的连续双射及 Y 到 X 上的连续双射；(ii) X 同胚于 Y 的一个开子空间，且 Y 同胚于 X 的一个开子空间。*
\end{enumerate}

*18) 设 X 为实平面 \(\mathbf{R}^2\)，设 R 为 X 上的等价关系，其等价类是诸直线 \(y = \beta\)（对 \(\beta > 0\)）以及诸半直线 \(x = \alpha\)，\(y \leqslant 0\)（对一切 \(\alpha \in \mathbb{R}\)）。证明：若把 X/R 视为 \(\mathfrak{P}_0(X)\) 的一个子集，则 X/R 上由 \(\mathcal{C}_{\Omega}\)（ \(\S 2\) ，习题 7）诱导的拓扑是离散的，而 X/R 上由 \(\mathcal{C}_{\Phi}\) 诱导的拓扑不是离散的，但使 X/R 成为非连通空间。最后证明：X/R 在商拓扑下是连通的。*

*¶ 19) 设 X 为局部紧空间。可以证明，X 可等同于在 § 9 习题 25 中构造的泛紧空间 \(\tilde{X}\) 的一个稠密开子空间（参见第九章，§ 1，习题 8）。再设 X 是连通且局部连通的。

\begin{enumerate}
\def\labelenumi{\alph{enumi})}
\item
  证明：对 X 的每个紧子集 K，X-K 的非相对紧的分支所成之集是有限的{[}参见习题 14 b){]}。由此推出：\(\tilde{X} - X\) 的一个点至多属于 X-K 的分支中一个的闭包{[}若 \(A_i\)（ \(i \leqslant i \leqslant n\)）是 X-K 的非紧分支，而 \(C_i\) 是 \(\tilde{X} - X\) 中位于 \(A_i\) 的闭包中的点所成的紧集，试构造一个以 X 与诸集 \(C_i\) 的和为底集的紧空间，使得 X 等同于该空间的一个稠密开子空间{]}。
\item
  设 \(R'\{x, y\}\) 为 \(\tilde{X}-X\) 上的下述等价关系：“对 X 的每个紧子集 K，x 与 y 位于 X-K 的同一个非紧分支的闭包中”。设 R 为 \(\tilde{X}\) 上的等价关系，其等价类是诸 mod \(R'\) 等价类以及由 X 的单个点组成的子集。证明：R 是豪斯多夫的，且 X 可等同于紧空间 \(X^{b}=\tilde{X}/R\) 的一个稠密开子空间；并证明：在紧空间 \(X^{b}-X\) 中，每个点都是它的既开又闭的邻域之交。\(X^{b}-X\) 的点称为空间 X 的端。
\item
  设 Y 为紧空间，使得 X 可等同于 Y 的一个稠密开子空间，且在紧空间 Y-X 中每个点都是它的既开又闭的邻域之交 (*)。证明：存在 \(X^{b}\) 到 Y 中的一个连续映射，它在 X 上诱导恒同映射。{[}若 Y-X 是两个不相交的非空闭集 A 与 B 的并，证明存在 Y 中两个不相交的开子集 U、V，使得 \(A \subset U\)、\(B \subset V\)，且使得 \(Y - (U \cup V)\) 是紧的{]}。
\item
  证明：若 X 是 σ-紧的，则 \(X^{b}-X\) 有一个可数基。
\item
  实直线有两个端，而 \(R^{n} (n \geqslant 2)\) 有一个。给出 \(R^{2}\) 的一个连通、局部连通、局部紧的子空间的例子，使其端所成之集具有连续统的势。*
\end{enumerate}

\begin{enumerate}
\def\labelenumi{\arabic{enumi})}
\setcounter{enumi}{19}
\tightlist
\item
  设 X 为豪斯多夫空间，其中每个点都有在 X 中既开又闭的邻域组成的基本系。设 Y 为 X 的子空间，A 为 Y 的一个在 Y 中既开又闭的紧子集；证明：存在 X 的一个在 X 中既开又闭的子集 B，使得 \(B \cap Y = A\)。
\end{enumerate}

¶ 21) a) 证明：拓扑空间 X 的下列条件等价：(i) 对 X 的每个开集 U，\(\overline{U}\) 是开的；(ii) 对 X 的每个闭集 F，\(\dot{F}\) 是闭的；(iii) 对 X 中每对不相交的开集 U、V，有 \(\overline{U} \cap \overline{V} = \varnothing\)。满足这些条件的豪斯多夫空间称为极不连通的；它这时是全不连通的。*有理直线是全不连通的但不是极不连通的。*

\begin{enumerate}
\def\labelenumi{\alph{enumi})}
\setcounter{enumi}{1}
\item
  豪斯多夫空间是极不连通的，当且仅当相伴随的半正则空间（§ 8，习题 20）是极不连通的。
\item
  设 X 为豪斯多夫空间，A 为 X 的稠密子集。证明：若 A 是极不连通的，且 X 的拓扑是 A-极大的（§ 3，习题 11），则 X 是极不连通的。
\item
  由 b) 与 c) 推出：若 X 是任一无限集，则 X 的（紧）超滤子空间是极不连通的（参见§ 9，习题 26 与 24）(**)。
\item
  设 X 为没有孤立点的豪斯多夫空间。证明：X 的拓扑是拟极大的（§ 2，习题 6），当且仅当 X 是次极大的（§ 8，习题 22）且极不连通的。（为证条件充分，利用 § 8 习题 22 e)，并证明 X 的每个没有孤立点的闭子集都是开的）。
\item
  证明：每个超正则空间（§ 8，习题 25）都是极不连通的，但没有孤立点的极不连通紧空间不是超正则的（参见§ 9，习题 27）。
\item
  证明：极不连通的半正则空间（§ 8，习题 20）是正则的（参见第二章，§ 4，习题 12）。
\item
  证明：在极不连通空间中，不存在具有无限多个不同项的收敛序列{[}用反证法，递归地构造两个两两不相交的开集序列 \((\mathbf{U}_{n})\)、\((\mathbf{V}_{n})\)，使得若 \(U = \bigcup_{n} U_{n}\) 且 \(V = \bigcup_{n} V_{n}\)，则 \(U \cap V = \varnothing\) 但 \(\overline{U} \cap \overline{V} \neq \varnothing\){]}。
\item
  证明：在极不连通空间 X 中，非孤立点 a 不能有同时既开又闭的、（关于关系 \(\supset\)）良序化的邻域基本系。（如 h) 中那样用反证法。）
\end{enumerate}

¶ 22) a) 证明：在极不连通空间中，每个开子空间与每个稠密子空间都是极不连通的。

\begin{enumerate}
\def\labelenumi{\alph{enumi})}
\setcounter{enumi}{1}
\item
  设 X 为豪斯多夫空间，其中每个点都有一个作为 X 的极不连通子空间的邻域。证明：X 是极不连通的。
\item
  给出一个豪斯多夫空间 X 的例子，它不是极不连通的，但有一个极不连通的稠密开子空间（把两个极不连通空间粘合起来）。
\item
  设 X 为可数离散空间，\(\tilde{X}\) 为 X 上的超滤子所成的（紧）空间（§ 9，习题 26）。设 \((\mathbf{A}_{n})_{n\in\mathbb{N}}\) 是把 X 分成无限子集的一个可数无限划分；在闭子空间 \(Y=\tilde{X}-X\)（\(\tilde{X}\) 的）中，诸集 \(B_{n}=\tilde{X}_{n}\cap Y\) 是开的、闭的且两两不相交的。设 \(B=\bigcup_{n}B_{n}\)；证明：B 在 Y 中的闭包 B 在 Y 中不是开的，从而 Y 不是极不连通的 (*)。（用反证法：利用习题 20，设 C 是 \(\tilde{X}\) 的一个既开又闭的子集，使得 \(C\cap Y=B\)；考察点 \(x_{n}\)（诸集 \(C\cap A_{n}\) 的每一个中的），设 \(\bar{J}\) 为诸点 \(x_{n}\) 所成之集 J 的闭包，并证明 \(\bar{J}\) 不与 B 相交）。
\end{enumerate}

\begin{enumerate}
\def\labelenumi{\arabic{enumi})}
\setcounter{enumi}{22}
\item
  给出双射映射 \(f\)（豪斯多夫空间 \(\mathbf{X}\) 到豪斯多夫空间 \(\mathbf{Y}\) 上的）的一个例子，使得在 \(f\) 下，\(\mathbf{X}\) 的每个紧子集的像是 \(\mathbf{Y}\) 的一个紧子集，且在 \(f\) 下，\(\mathbf{X}\) 的每个连通子集的像是 \(\mathbf{Y}\) 的一个连通子集，但 \(f\) 不是连续的（参见§ 9，习题 4）。
\item
  设 X 为拓扑空间。
\end{enumerate}

\begin{enumerate}
\def\labelenumi{\alph{enumi})}
\tightlist
\item
  设集 \(\mathfrak{P}_{\Omega}(\mathbf{X})\) 赋有诸拓扑 \(\mathfrak{C}_{\Omega}, \mathfrak{C}_{\Phi}\)（ \(\S 2\) ，习题 7）或 \(\mathfrak{C}_{\Theta}\)（ \(\S 8\) ，习题 12）之一。证明：若子集 \(\mathfrak{B}\)（\(\mathfrak{P}_{\Omega}(\mathbf{X})\) 的）是连
\end{enumerate}

通的，且每个 \(\mathbf{M} \in \mathfrak{B}\) 都是连通的，则 \(\bigcup_{\mathbf{M} \in \mathfrak{M}} \mathbf{M}\) 是连通的。

\begin{enumerate}
\def\labelenumi{\alph{enumi})}
\setcounter{enumi}{1}
\tightlist
\item
  设 \(\mathfrak{S}\) 为 \(\mathfrak{P}_{0}(\mathbf{X})\) 的一个子集，它包含 \(\mathbf{X}\) 的所有有限非空子集所成之集。证明：若 \(\mathbf{X}\) 是连通的，则 \(\mathfrak{S}\) 也是连通的{[}利用 § 2 习题 7 b) 与 § 4 命题 8，并考察诸映射 \((x_{1},\ldots ,x_{n})\to \{x_{1},\ldots ,x_{n}\} ]\)。
\end{enumerate}

\begin{enumerate}
\def\labelenumi{\arabic{enumi})}
\setcounter{enumi}{24}
\tightlist
\item
  给定两个拓扑空间 X、Y，映射 \(f: \mathbf{X} \to \mathbf{Y}\) 称为局部同胚，如果对每个 \(x \in \mathbf{X}\) 存在 \(x\) 的一个邻域 U，使得 \(f(\mathbf{U})\) 是 \(f(x)\) 在 Y 中的一个邻域，且映射 \(\mathbf{U} \to f(\mathbf{U})\)（在 U 上与 \(f\) 相合）是一个同胚。每个局部同胚都是连续的开映射。
\end{enumerate}

\begin{enumerate}
\def\labelenumi{\alph{enumi})}
\item
  设 X 是豪斯多夫的，Y 是局部连通的。设 \(f: \mathbf{X} \to \mathbf{Y}\) 为局部同胚，具有这样的性质：存在整数 \(n > 0\)，使得 \(\overline{f}^1(y)\) 恰有 \(n\) 个点（对每个 \(y \in \mathbf{Y}\)）。证明：每个 \(y \in \mathbf{Y}\) 都有一个开邻域 V，使得 \(\overline{f}^1(V)\) 有 \(n\) 个分支 \(\mathbf{U}_i\)（ \(1 \leqslant i \leqslant n\)），且映射 \(\mathbf{U}_i \to \mathbf{V}\)（它与 \(f\) 在 \(\mathbf{U}_i\) 上相合）是同胚。这时 \(f\) 是逆紧映射。
\item
  设 X 是豪斯多夫的，且 Y 是连通且局部连通的。设 \(f: \mathbf{X} \to \mathbf{Y}\) 为逆紧的局部同胚。证明：\(f\) 满足 \(a\) ) 的假设。{[}若对每个整数 \(n > 0\)，\(\mathbf{Y}_n\) 是这样的 \(y \in \mathbf{Y}\) 所成之集：\(\bar{f}^1(y)\) 至少有 \(n\) 个点，证明 \(\mathbf{Y}_n\) 在 \(\mathbf{Y}\) 中既开又闭{]}。
\end{enumerate}

\(^{*}c)\) 给出满射局部同胚 \(f: X \to Y\) 的一个例子，其中 Y 是紧的、连通的且局部连通的，X 是局部紧的、连通的且局部连通的，\(\overline{f}^{1}(y)\) 对每个 \(y \in Y\) 至多由两个点组成，但 f 不是逆紧的（参见§ 5，习题 4）。
% semantic-fragments: legacy-input-seam
\relax{}
